\documentclass[10pt,a4paper]{amsart}
\usepackage[margin=19mm]{geometry}
\usepackage{amsmath,amssymb,mathtools}
\usepackage[scr=rsfs,cal=euler]{mathalfa}
\usepackage{xfrac}
\usepackage[unicode,hidelinks,bookmarks=false]{hyperref}
\newcommand{\ie}{i.e.\ }
\newcommand{\cf}{cf.\ }
\newcommand{\resp}{respectively\ }
\newcommand{\Subsection}{\subsection}
\providecommand{\nobreakdash}{\nobreak\hbox{-}}
\setlength{\emergencystretch}{2em}
\multlinegap0pt
\usepackage{csquotes}
\usepackage{xcolor}
\usepackage{amssymb}
\usepackage[normalem]{ulem}
\newtheorem{prop}{Proposition}
\title{On Regular Higher Power Rational Diophantine Triples}
\author{Alen Andra\v{s}ek}
\date{Source: arXiv:2604.17018v2 (CC BY 4.0)}
\begin{document}
\maketitle

\noindent\emph{Source and licence.} On Regular Higher Power Rational Diophantine Triples. Version arXiv:2604.17018v2 (CC BY 4.0), distributed by arXiv under the Creative Commons Attribution 4.0 International licence, http://creativecommons.org/licenses/by/4.0/, as recorded in the arXiv licence metadata for this exact version and shown on the abstract page https://arxiv.org/abs/2604.17018v2. Full licence notice: \texttt{licenses/arxiv-2604.17018-CC-BY-4.0.txt}.\par\smallskip

\noindent\textit{Editorial scope.} Counted unit: the article's prop equation\_Weierstrass\_form (source0.tex lines 274--280). The article's complete original proof of it is delivered with it (source0.tex lines 281--304, one \texttt{proof} environment of 24 lines). No mathematics is written, repaired or re-derived: the statement and the proof are the article's own lines, in the article's own notation, and no retained line is substituted or re-flowed.  Retained uncounted as the source-local context the proof needs: lines 270--272.  Nothing was omitted from any retained span, so the package carries no omission record.  No retained line contains a citation, so the wrapper carries no bibliography and no bibliography entry.  No retained line contains a figure environment, an image-inclusion command or drawing code, so no image is delivered and the package declares no asset dependency.  Editorial formatting only, in addition: an AMS \texttt{amsart} preamble that restates the article's own declarations and macros used by the retained lines, namely the environments \texttt{prop} on separate counters, together with the standard packages those lines need.  The retained environments print in this document as Proposition 1 in order of appearance, whereas the article prints the same items with the numbering of its own full document.  The complete native source of the article as distributed in the arXiv e-print for this exact version is bundled unchanged as \texttt{sources/198714/source0.tex}.
\begin{align} \label{s^2-r^2 equation_quartic_form}
    (u^2+1)^2r^4 - (u^2-1)^2 &= \square.
\end{align}

\begin{prop}
    The curve (\ref{s^2-r^2 equation_quartic_form}) has the Weierstrass form \begin{align} \label{s^2-r^2 equation_Weierstrass_form}
        y^2 = x^3 + 4(u^4-1)^2x.
    \end{align}
    This curve has torsion $\mathbb{Z}/2\mathbb{Z} = \{\mathcal{O}, (0,0)\}$ and a point of infinite order $P = (2(u^2+1)(u-1)^2, 4(u^2+1)^2(u-1)^2)$ over $\mathbb{Q}(u)$.\\
    A point on this curve gives a regular triple with $(r,s)=\left(\frac{-y}{2x(u^2+1)}, \frac{-y}{2x(u^2-1)}\right)$.
\end{prop}

\begin{proof} Rewrite (\ref{s^2-r^2 equation_quartic_form}) as
\begin{align*}
    r^4-\alpha^2 = w^2
\end{align*}
where $\alpha = \frac{u^2-1}{u^2+1}$, for some rational $u$. Put $w = W + r^2$ to get
\begin{align*}
    -\alpha^2 = W^2+2Wr^2\Big /\cdot W
\end{align*}
Now, set $z = Wr$ to arrive at the equation
\begin{align*}
    W^3+\alpha^2W = -2z^2
\end{align*}
Finally, multiplying by $8$, and setting $-2W = X$, $4z=Z$, we get the Weierstrass equation
\begin{align*}
    Z^2 =  X^3+4\alpha^2 X.
\end{align*}
Plugging in $\alpha=\frac{u^2-1}{u^2+1}$ and setting $x = (u^2+1)^2 X$ and $y = (u^2+1)^3Z$, we finally get \[y^2=x^3+4(u^4-1)^2x.\]
Combining the transformation formulas, we find that
\begin{align*}
    r = \frac{-y}{2x(u^2+1)}. 
\end{align*}
Then $s=r \frac{u^2+1}{u^2-1} = \frac{-y}{2x(u^2-1)}$. 
% (The minus sign may as well be omitted since $r,s$ appear only squared in \ref{s^2-r^2 equation_quartic_form}.)
\end{proof}

\end{document}
